% Generated draft: requires mathematical and source review.
\documentclass[11pt]{article}
\usepackage[margin=25mm]{geometry}
\usepackage{fontspec}
\setmainfont{Latin Modern Roman}
\usepackage{amsmath,amssymb,mathtools,unicode-math}
\setmathfont{Latin Modern Math}
\usepackage{enumitem,needspace,xurl,hyperref}
\setlength{\parindent}{0pt}
\setlength{\parskip}{5pt}
\setlength{\emergencystretch}{3em}
\newcommand{\N}{\mathbb{N}}
\newcommand{\Z}{\mathbb{Z}}
\newcommand{\Q}{\mathbb{Q}}
\newcommand{\R}{\mathbb{R}}
\newcommand{\C}{\mathbb{C}}
\newcommand{\F}{\mathbb{F}}
\newcommand{\A}{\mathbb{A}}
\newcommand{\PP}{\mathbb{P}}
\newcommand{\E}{\mathbb{E}}
\newcommand{\eps}{\varepsilon}
\newcommand{\dd}{\,\mathrm{d}}
\begin{document}
\section*{UnsolvedMath ID: 2000002}
\textbf{Source problem code:} catalogue entry
\subsection{Definitions and mathematical statement}
A stored-energy density \(W\) in nonlinear elasticity assigns an energy to a deformation gradient \(F\). Anisotropy means that material directions can affect the response; the source does not specify a material symmetry group or an explicit family of energies.

For square matrices \(F\in\R^{d\times d}\), let \(\mathcal M(F)\) be the vector of all minors of orders \(1,\ldots,d\). A conventional polyconvexity condition is the existence of a convex function \(P\) on the minors space such that
\[
W(F)=P(\mathcal M(F)).
\]
For a specified energy class in which the integrals are defined, a conventional quasiconvexity condition is
\[
|\Omega|W(F)\le\int_\Omega W(F+\nabla\varphi(x))\,\dd x
\]
for every matrix \(F\), every bounded Lipschitz domain \(\Omega\subset\R^d\), and every \(\varphi\in W^{1,\infty}_0(\Omega;\R^d)\). Here \(W^{1,\infty}_0\) means Lipschitz perturbations with zero boundary values. Energies defined only for \(\det F>0\) are often treated by extension to \(+\infty\) outside that set; the admissible domain and treatment of such singular energies must be fixed when applying either definition.

Ball's request is to develop useful criteria or methods for deciding or proving these two conditions for stored-energy functions arising in anisotropic nonlinear elasticity. It is a methodological research problem, not a specified claim that every anisotropic energy satisfies either condition. No particular dimension, constitutive model, regularity, growth assumptions, material symmetries, meaning of usefulness, or completeness requirement for a test is supplied. The definitions above give the standard mathematical properties to be tested, conditional on an appropriate chosen energy class; they do not fill in an absent model.

The supplied OPEN-TRIAGE summary reports substantial positive polyconvexity results for many anisotropic models and only partial/numerical tools for general quasiconvexity verification. It does not supply particular models, theorem hypotheses, or supporting references, and does not claim a complete general solution.
\subsection{Short English statement}
Develop usable tests for polyconvexity and quasiconvexity of anisotropic elastic stored energies. The record specifies the research aim but no model class or exact criterion of a successful verification method.
\subsection{Sources}
\begin{itemize}
\item Local UnsolvedMath registry, permanent ID \(2000002\), exact source code \texttt{catalogue entry}, ``Some Open Problems in Elasticity --- Testing convexity conditions.''
\item Ball, elasticity problem list (2002), Problem 2; supplied source \url{https://people.maths.ox.ac.uk/ball/Articles%20in%20Conference%20Proceedings%20and%20Books/JMB%202002%20re%20Marsden%2060th.pdf}. PDF-text extraction; recorded access 2026-07-29.
\item Supplied OPEN-TRIAGE summary reports major partial progress but gives no specific authoritative references for the described verification toolkit or current limits.
\end{itemize}
\end{document}
